\documentclass[11pt]{article}
\usepackage[margin=1in]{geometry}
\usepackage{amsmath}
\usepackage{graphicx}
\usepackage{booktabs}
\usepackage[hidelinks]{hyperref}

\title{Adversarial Fine-Tuning and Robust Overfitting in V-JEPA 2}
\author{Kaiwen Du}
\date{Draft research protocol --- October 2026}

\begin{document}
\maketitle

\begin{abstract}
This document outlines a planned study of adversarial fine-tuning of
V-JEPA 2 under its native masked latent prediction loss. The target encoder
will remain fixed while the context encoder and predictor are updated.
We will examine training and held-out robustness over time, validate the
attacks, and measure downstream action-recognition consequences.
No experiments or results are reported in this draft.
\end{abstract}

\section{Introduction}
V-JEPA 2 learns video representations through latent prediction
\cite{assran2025vjepa2}. Robust overfitting has been observed in
adversarially trained classifiers \cite{rice2020overfitting}.
This project asks whether a related pattern occurs during fixed-target
adversarial adaptation of a pretrained latent video predictor.

\section{Related Work}
The dated literature audit in the repository records nearby work on
representation attacks, adversarial self-supervised learning, and video
robustness. Expand this section from verified primary sources before
submission. No priority claim is made in this draft.

\section{Proposed Method}
Let $f_{\theta}$ be the context encoder, $g_{\phi}$ the predictor,
and $f_{\bar{\theta}}$ the fixed target encoder. For a video $x$ and
mask $M$, define the native masked prediction loss as
\begin{equation}
  \mathcal{L}(x,\delta,M) =
  \operatorname{mean}\left|
  g_{\phi}(f_{\theta}(x+\delta;M);M)
  - \operatorname{sg}\!\left[T_M(f_{\bar{\theta}}(x))\right]
  \right|,
\end{equation}
where $T_M$ applies upstream target normalization and target-token
selection, and $\operatorname{sg}$ stops gradients. The adversary
maximizes this loss subject to $\|\delta\|_{\infty}\leq\epsilon$
and valid pixel bounds. The outer update minimizes the attacked loss
over $\theta$ and $\phi$; $\bar{\theta}$ stays fixed.

\section{Experimental Protocol}
The authoritative draft settings are in \texttt{configs/study.yaml},
with methodology in \texttt{docs/experimental\_protocol.md}.
The study compares a pretrained reference, clean fine-tuning, and
adversarial fine-tuning on group-separated UCF101 subsets. Three
training seeds are planned for each fine-tuning condition.

\section{Results}
\textbf{No results available.} Insert only measurements from traceable
runs after the implementation and attack checks pass. Report individual
seeds, checkpoint selection, uncertainty, and negative findings.

\section{Limitations}
The target encoder is fixed, so this study does not establish the
behavior of the original EMA pretraining procedure. A small adaptation
dataset and three seeds limit generalization. Pretraining overlap cannot
be ruled out. Masked video prediction is not itself a demonstration of
causal future prediction or reliable robot planning.

\section{Reproducibility}
The final report will identify code and upstream revisions, model and
manifest hashes, resolved configurations, seeds, attack settings,
hardware, environment, and compute cost. No training-ready environment
or experiment artifact release exists at this stage.

\bibliographystyle{plain}
\bibliography{references}
\end{document}
